\chapter{Childhood illnesses}
\index{Childhood illnesses}

\medskip
\noindent\textit{\textbf{Under review.} This chapter is an AI-assisted educational draft; it carries no source citations and is not a substitute for professional medical advice.}

\section*{Definition and Overview}
Childhood illnesses are the infections and other conditions that affect children, particularly the common, mostly self-limiting infections of the first years of life. These include colds, influenza, gastroenteritis, chickenpox, hand, foot and mouth disease, otitis media (middle-ear infection), croup, and bronchiolitis, as well as common non-infectious problems such as eczema, asthma, and allergies. Most childhood illnesses are caused by viruses, run a predictable course, and resolve on their own with supportive care. A small number are caused by bacteria and need antibiotics, and a few serious infections, such as meningitis, require urgent treatment.

\section*{Epidemiology}
Young children get more infections than any other age group. It is normal for a child in childcare or with older siblings to have 6 to 10 viral infections a year, especially in the first 2 to 3 years of life as their immune systems meet new viruses. Most are self-limiting. Common examples include bronchiolitis, which affects almost all children in the first 2 years; otitis media, which follows many colds; and gastroenteritis, one of the most common reasons for seeing a doctor in early childhood. Immunisation has made serious illnesses such as measles, whooping cough, and meningitis much rarer.

\section*{Aetiology and Pathophysiology}
Children are vulnerable to frequent infections for several reasons. Their immune systems have not yet encountered most viruses, so every new exposure is a first infection. They also have smaller airways and shorter, more horizontal eustachian tubes, which is why croup, bronchiolitis, and ear infections are so common. Viruses spread easily through childcare and schools via hands, coughing, and sneezing. Most childhood infections involve viruses:

\begin{itemize}
    \item \textbf{Respiratory viruses:} cold viruses, influenza, respiratory syncytial virus (RSV) causing bronchiolitis
    \item \textbf{Gut viruses:} rotavirus and norovirus causing gastroenteritis
    \item \textbf{Rash viruses:} chickenpox, hand, foot and mouth disease, roseola
    \item \textbf{Bacterial infections:} ear infections, tonsillitis, and urinary tract infections, sometimes following a viral illness
\end{itemize}
\section*{Clinical Features}
Symptoms depend on the illness but commonly include:

\begin{itemize}
    \item Fever, which is the most common sign and usually peaks in the afternoon or evening
    \item Runny or blocked nose, cough, and a sore throat
    \item Earache or tugging at the ears in young children
    \item Vomiting and diarrhoea with gastroenteritis
    \item A hoarse, barking cough with croup, often worse at night
    \item Fast or noisy breathing with bronchiolitis in babies
    \item A rash, which may be spotty, blistery, or blotchy
    \item Reduced appetite, tiredness, and clinginess
\end{itemize}
Most symptoms settle within a few days to a week, though a cough can linger for several weeks.

\section*{Differential Diagnosis}
\begin{itemize}
    \item Viral vs.\ bacterial infection -- the commonest clinical question, answered mainly by the pattern of illness and examination
    \item Croup vs.\ other causes of a barking cough, such as a foreign object inhaled into the airway
    \item Bronchiolitis vs.\ asthma vs.\ pneumonia in a wheezy infant
    \item A viral rash vs.\ a serious illness such as measles or meningococcal disease
    \item A fever with no clear source, which may need investigation to exclude a urinary infection or meningitis
\end{itemize}

\section*{When to Seek Medical Care}
See a doctor if a child has a fever lasting more than a few days, symptoms that are getting worse, an earache, trouble drinking, or any rash you are unsure about. Babies under 3 months with a fever should always be seen.

\textbf{Contact your local emergency services or seek emergency care if your child:}
\begin{itemize}
    \item Is drowsy, floppy, or hard to wake
    \item Breathes rapidly, with sucking in of the ribs, or struggles to breathe
    \item Has a fit or seizure
    \item Develops a rash of small purple or red spots that do not fade when pressed (possible meningococcal disease)
    \item Cannot keep any fluids down and shows signs of dehydration (dry mouth, sunken eyes, little urine)
\end{itemize}

\section*{Diagnosis and Investigations}
\begin{itemize}
    \item \textbf{History and examination:} the pattern and timing of symptoms, and checking the throat, ears, chest, and skin
    \item \textbf{Observation:} how the child looks, behaves, and interacts is often the best guide to how unwell they are
    \item \textbf{Tests when needed:} throat swab, urine test, or blood tests if a bacterial infection is suspected
    \item \textbf{Chest X-ray:} where pneumonia is suspected
    \item \textbf{Oximetry:} a probe measuring oxygen levels in babies with breathing difficulty
\end{itemize}

\section*{Management}
\begin{itemize}
    \item \textbf{Supportive care:} rest, fluids, and paracetamol or ibuprofen at weight-appropriate doses for fever and pain
    \item \textbf{Antibiotics:} only for bacterial infections such as some ear infections, tonsillitis, or pneumonia -- they do not work for viruses
    \item \textbf{Oral rehydration solution:} for gastroenteritis, to replace lost fluids and salts
    \item \textbf{Steam or cool air:} for croup and blocked noses in older children
    \item \textbf{Hospital care:} for severe bronchiolitis, dehydration, or any child who needs oxygen or close monitoring
    \item \textbf{Home comfort:} small, frequent drinks, and keeping the child cool but not overdressed
\end{itemize}

\section*{Complications}
\begin{itemize}
    \item Dehydration, especially with gastroenteritis in young children
    \item Ear infections that recur or cause a persistent middle-ear effusion
    \item Pneumonia or breathing difficulty complicating a chest infection
    \item Febrile convulsions with a rapid temperature rise -- frightening but usually harmless
    \item Rarely, a viral illness such as influenza can be severe and require hospital care
    \item Untreated serious bacterial infection (meningitis, sepsis) can progress rapidly
\end{itemize}

\section*{Prognosis and Outcomes}
Most childhood illnesses are mild and resolve on their own within a week or so. Children recover quickly and rarely have lasting effects. A cough may persist for several weeks after a viral infection and is usually not a cause for concern. Serious complications are uncommon in otherwise healthy, immunised children. Frequent infections in the early years usually settle as the child's immune system matures and their exposure to new viruses falls once they have met them.

\section*{Prevention and Self-Care}
\begin{itemize}
    \item Keep vaccinations up to date on the immunisation schedule
    \item Encourage regular hand washing, especially before meals and after coughing or sneezing
    \item Teach children to cover coughs and sneezes with the elbow
    \item Keep a sick child home from childcare or school until they are well enough to return
    \item Do not give aspirin to children; use paracetamol or ibuprofen at the right dose if needed
    \item Offer plenty of fluids and let the child rest
    \item Seek medical advice early for babies, or for any child who is getting worse
\end{itemize}

\section*{Sources}
\begin{itemize}
    \item World Health Organization immunization guidance: \url{https://www.who.int/health-topics/vaccines-and-immunization}
    \item World Health Organization guidance on integrated management of childhood illness: \url{https://www.who.int/teams/maternal-newborn-child-adolescent-health-and-ageing/child-health/integrated-management-of-childhood-illness}
\end{itemize}
